% Generated draft: requires mathematical and source review.
\documentclass[11pt]{article}
\usepackage[margin=25mm]{geometry}
\usepackage{fontspec}
\setmainfont{Latin Modern Roman}
\usepackage{amsmath,amssymb,mathtools,unicode-math}
\setmathfont{Latin Modern Math}
\usepackage{enumitem,needspace,xurl,hyperref}
\setlength{\parindent}{0pt}
\setlength{\parskip}{5pt}
\setlength{\emergencystretch}{3em}
\newcommand{\N}{\mathbb{N}}
\newcommand{\Z}{\mathbb{Z}}
\newcommand{\Q}{\mathbb{Q}}
\newcommand{\R}{\mathbb{R}}
\newcommand{\C}{\mathbb{C}}
\newcommand{\F}{\mathbb{F}}
\newcommand{\A}{\mathbb{A}}
\newcommand{\PP}{\mathbb{P}}
\newcommand{\E}{\mathbb{E}}
\newcommand{\eps}{\varepsilon}
\newcommand{\dd}{\,\mathrm{d}}
\begin{document}
\section*{UnsolvedMath ID: 30006578}
\textbf{Source problem code:} catalogue entry
\subsection{Definitions and mathematical statement}
Let
\[
D=\{z\in\C:|z|\leq 1\},\qquad D^n=\{(z_1,\ldots,z_n)\in\C^n:|z_i|\leq 1\text{ for every }i\}.
\]
The supplied statement denotes by
\[
C_\bullet=\widetilde{\Omega}_{\infty-\bullet}^{\mathrm{alg}}(D_\infty)
\]
a boundary-vanishing complex of codegree forms whose coefficients are holomorphic on neighborhoods of the closed polydiscs and algebraic over \(\Q(z_1,z_2,\ldots)\). Here the variables \(z_i\) are the complex coordinate variables. The supplied record does not specify the construction of \(D_\infty\), the codegree terms, the differential, the boundary restrictions, or the transition maps between finite-dimensional stages. Consequently, this notation identifies the intended source-dependent complex but does not give a complete independent definition of it.

\textbf{Conjectural target (Conjecture 2 in the supplied excerpt).} With the intended source definitions, the complex is acyclic outside degree zero:
\[
H_j(C_\bullet)=0\qquad\text{for every degree }j\ne 0.
\]
Moreover, the formal integration map
\[
\mathcal I:H_0(C_\bullet)\longrightarrow\C,
\qquad
\mathcal I([f\,\mathrm dz])=\int_{[0,1]^\infty}f\,\mathrm dz,
\]
is injective; equivalently, for every \(\xi\in H_0(C_\bullet)\),
\[
\mathcal I(\xi)=0\quad\Longrightarrow\quad\xi=0.
\]
The displayed integral reproduces the formal notation in the supplied excerpt. Its meaning, the notation \(\mathrm dz\), and its descent to homology require the source construction, including the unavailable equation (1); no literal infinite-dimensional integration measure is specified here.

The supplied Remark 3 states that the image of \(\mathcal I\) is precisely the ring of effective periods, that injectivity is essentially a reformulation of the Kontsevich--Zagier period conjecture, and that this identification is nontrivial and relies on reference [5]. It identifies the period conjecture with reference [6] and relates the acyclicity assertion closely to the existence of a motivic \(t\)-structure. These are supplied source qualifications, not independently verified equivalences or additional proved assertions.
\subsection{Short English statement}
For the intended algebraic complex of boundary-vanishing codegree forms on closed polydiscs, the conjecture asks that homology vanish in every nonzero degree and that formal integration distinguish all degree-zero homology classes. The supplied record omits definitions needed to specify the complex and integration map independently. Its open status and its connection with the Kontsevich--Zagier period conjecture are supplied assessments requiring review.
\subsection{Sources}
\begin{itemize}
\item Local UnsolvedMath registry: supplied record with permanent ID \(30006578\), title \emph{Algebraic Effective-Period Complex Conjecture}, exact source problem code \texttt{catalogue entry}, in the Oberwolfach Reports Open Problems collection. The record assigns proposed year 2026 and supplies no named proposer.
\item \emph{Model Theory and Arithmetic} (2026), Oberwolfach Reports, DOI 10.4171/owr/2026/3, pp.\ 157--216. Supplied source and literature link, labeled \emph{Oberwolfach Report 3/2026}: \url{https://doi.org/10.4171/owr/2026/3}. The excerpt identifies Conjecture 2 and Remark 3, but supplies no exact page within the report.
\item The supplied literature assessment was checked on \texttt{2026-08-22} and labels the problem open, connecting acyclicity and period-map injectivity with the conjectural framework described above. This is supplied dated curation derived from the OWR record, not an independently verified literature assessment. The supplied background expressly requests independent verification of time-sensitive status and cited claims.
\item The registry labels the statement \texttt{corrected\_verified} and reports that the official report was checked directly to repair OCR and notation. That verification claim is part of the supplied record; it has not been independently confirmed here.
\item The excerpt refers to equation (1) and references [5] and [6] without supplying their contents or bibliographic entries. It also contains a reference to [1], Remarque 2.109 and Conjecture A, whose attachment to this conjecture or the following adic-period discussion is unclear from the extraction.
\end{itemize}
\end{document}
